\paragrafos{Δυνάμεις φυσικών αριθμών: ποιος πάει πρώτος}{Α.1.3, «Δυνάμεις φυσικών αριθμών»}

\begin{anagke}
Παρασκευή απόγευμα. Η Λίνα φτιάχνει κωδικό για το καινούργιο της κινητό: τέσσερα ψηφία. «Πόσοι διαφορετικοί κωδικοί υπάρχουν;» ρωτάει η Ζωή. Η Λίνα σκέφτεται δυνατά: «Για το πρώτο ψηφίο έχω 10 επιλογές, από το 0 ως το 9. Για το δεύτερο πάλι 10, και για το τρίτο, και για το τέταρτο». Ο Ηλίας γράφει $10\cdot 10\cdot 10\cdot 10$ και βρίσκει 10.000. «Και αν ο κωδικός είχε είκοσι ψηφία;» ρωτάει. Γράφει είκοσι δεκάρια στη σειρά και χάνει το μέτρημα στη μέση.

Ο Άρης, που βιάζεται πάντα, γράφει σε ένα χαρτάκι $2+3\cdot 4$ και ρωτάει τη Ζωή πόσο κάνει. «Δεκατέσσερα», λέει εκείνη. «Είκοσι!» φωνάζει ο Άρης. «Δύο και τρία πέντε, επί τέσσερα, είκοσι». Ο Θοδωρής, που περνάει από δίπλα, τους κοιτάει και λέει μόνο: «Και οι δύο έχετε ένα κομπιουτεράκι στην τσέπη. Δείτε τι βγάζει». Το κομπιουτεράκι δίνει 14.

Δύο προβλήματα. Το πρώτο: πώς γράφεται σύντομα ο ίδιος παράγοντας πολλές φορές; Το δεύτερο: όταν μια παράσταση έχει πολλές πράξεις, ποια γίνεται πρώτη, και ποιος το αποφασίζει;
\end{anagke}

\begin{lexi}{δύναμη, εκθέτης, βάση}{δύναμη, εκθέτης}
Η λέξη \lexiword{δύναμη} στα μαθηματικά δεν έχει σχέση με τα μπράτσα. Ο αριθμός «έχει τη δύναμη» να πολλαπλασιάζεται με τον εαυτό του πολλές φορές. Το $2^{10}$ διαβάζεται «το δύο στη δέκατη».

Ο \lexiword{εκθέτης} είναι το μικρό νούμερο πάνω δεξιά. Το «εκθέτω» σημαίνει «δείχνω»: ο εκθέτης δείχνει πόσες φορές μπαίνει ο παράγοντας. Ο αριθμός που πολλαπλασιάζεται λέγεται \lexiword{βάση}, γιατί πάνω του «στηρίζεται» το σύμβολο.
\end{lexi}

\begin{metafrasi}
Το 2 πολλαπλασιασμένο με τον εαυτό του 10 φορές. & $2^{10}=2\cdot 2\cdot\dots\cdot 2$ (10 παράγοντες)\\
Το τετράγωνο του 5. & $5^2=5\cdot 5=25$\\
Ο κύβος του 3. & $3^3=3\cdot 3\cdot 3=27$\\
Το 7 στην πρώτη. & $7^1=7$\\
Ένα εκατομμύριο, με δύναμη του 10. & $10^6=1.000.000$\\
Πρώτα οι δυνάμεις, μετά οι πολλαπλασιασμοί. & $2+3\cdot 4^2=2+3\cdot 16=50$\\
\end{metafrasi}

\tmima{Η θεωρία}

\begin{orismos}{δύναμη}
Το γινόμενο $\alpha\cdot\alpha\cdot\alpha\cdots\alpha$ με $\nu$ παράγοντες ίσους με το $\alpha$ λέγεται \textbf{δύναμη} του $\alpha$ στη $\nu$ και γράφεται $\alpha^{\nu}$. Ο $\alpha$ είναι η \textbf{βάση} και ο $\nu$ ο \textbf{εκθέτης}.
\end{orismos}

Η δεύτερη δύναμη $\alpha^2$ λέγεται και \textbf{τετράγωνο} του $\alpha$ (είναι το εμβαδό τετραγώνου με πλευρά $\alpha$). Η τρίτη δύναμη $\alpha^3$ λέγεται και \textbf{κύβος} του $\alpha$ (είναι ο όγκος κύβου με ακμή $\alpha$). Η πρώτη δύναμη $\alpha^1$ είναι ο ίδιος ο αριθμός $\alpha$. Οι δυνάμεις του 1 είναι όλες 1: $1^{\nu}=1$.

\begin{prosoxi}
Το $2^3$ δεν είναι $2\cdot 3=6$ αλλά $2\cdot 2\cdot 2=8$. Ο εκθέτης δείχνει πόσες φορές γράφεται η βάση, όχι επί πόσο την πολλαπλασιάζουμε.
\end{prosoxi}

\tmimasmall{Οι δυνάμεις του 10}

Οι δυνάμεις του 10 έχουν τόσα μηδενικά όσος ο εκθέτης:
$$10^2=100,\qquad 10^3=1.000,\qquad 10^6=1.000.000.$$
Γι' αυτό η δεκαδική μας γραφή αναλύεται εύκολα. Για παράδειγμα:
$$7.604=7\cdot 10^3+6\cdot 10^2+0\cdot 10+4.$$

\tmimasmall{Αριθμητική παράσταση και προτεραιότητα}

Μια \textbf{αριθμητική παράσταση} είναι μια σειρά αριθμών που συνδέονται με σύμβολα πράξεων. Ο αριθμός που βγαίνει όταν κάνουμε όλες τις πράξεις λέγεται \textbf{τιμή} της. Για να βγάζουμε όλοι την ίδια τιμή, όλος ο κόσμος έχει συμφωνήσει σε μια σειρά.

\begin{kanonas}{προτεραιότητα των πράξεων}
\begin{enumerate}[leftmargin=1.6em,itemsep=0.2ex,topsep=2pt]
\item Πρώτα οι πράξεις μέσα στις \textbf{παρενθέσεις}, με την ίδια σειρά που θα βρεις παρακάτω.
\item Μετά οι \textbf{δυνάμεις}.
\item Μετά οι \textbf{πολλαπλασιασμοί και οι διαιρέσεις}, από αριστερά προς τα δεξιά.
\item Τελευταίες οι \textbf{προσθέσεις και οι αφαιρέσεις}, από αριστερά προς τα δεξιά.
\end{enumerate}
\end{kanonas}

Η συμφωνία αυτή δεν είναι «φυσικός νόμος». Είναι συμφωνία, όπως ο κανόνας ότι στον δρόμο πάμε δεξιά. Αν δεν υπήρχε, το $2+3\cdot 4$ θα είχε δύο απαντήσεις (14 και 20) και δύο μηχανικοί δεν θα συμφωνούσαν ποτέ σε έναν υπολογισμό. Με τη συμφωνία η απάντηση είναι μία, 14, και η Ζωή έχει δίκιο. Αν θέλουμε το 20, γράφουμε $(2+3)\cdot 4$: η παρένθεση λέει ρητά «αυτό πρώτα».

\begin{prosoxi}
\begin{itemize}[leftmargin=1.2em,itemsep=0.3ex,topsep=2pt]
\item Στο $2\cdot 3^2$ ο εκθέτης αφορά μόνο το 3: $2\cdot 9=18$. Για να τετραγωνίσεις όλο το γινόμενο, γράφεις $(2\cdot 3)^2=36$.
\item Οι πράξεις του ίδιου επιπέδου γίνονται από αριστερά προς τα δεξιά: $24:4:2=6:2=3$, όχι $24:(4:2)=12$.
\end{itemize}
\end{prosoxi}


\tmimasmall{Ιστορική αναδρομή: πώς έγραφαν τους αριθμούς}

Το δεκαδικό σύστημα, που σου φαίνεται αυτονόητο, χρειάστηκε χιλιάδες χρόνια για να διαμορφωθεί. Οι Σουμέριοι και οι Βαβυλώνιοι (3η–2η χιλιετία π.Χ.) έγραφαν τους αριθμούς σε πήλινες πινακίδες. Οι Αιγύπτιοι είχαν ένα ιερογλυφικό για κάθε τάξη: 1, 10, 100, 1.000, \dots, 1.000.000. Οι Ίωνες, τον 5ο αιώνα π.Χ., έφτιαξαν το \textbf{αλφαβητικό σύστημα}: κάθε γράμμα του αλφαβήτου σήμαινε έναν αριθμό.

\begin{center}
\small\renewcommand{\arraystretch}{1.25}\setlength{\tabcolsep}{4pt}
\begin{tabular}{cccccccccc}
α & β & γ & δ & ε & \arch{ϛ} & ζ & η & θ & \\
1 & 2 & 3 & 4 & 5 & 6 & 7 & 8 & 9 & \\
ι & κ & λ & μ & ν & ξ & ο & π & \arch{ϙ} & \\
10 & 20 & 30 & 40 & 50 & 60 & 70 & 80 & 90 & \\
ρ & σ & τ & υ & φ & χ & ψ & ω & \arch{ϡ} & \\
100 & 200 & 300 & 400 & 500 & 600 & 700 & 800 & 900 & \\
\end{tabular}
\end{center}

Έτσι το 831 γράφεται ωλα ($800+30+1$). Οι \textbf{Ρωμαίοι} είχαν σύμβολα για τους αριθμούς $1,\,5,\,10,\,50,\,100,\,500,\,1.000$:

\begin{center}
\small\setlength{\tabcolsep}{6pt}
\begin{tabular}{ccccccc}
I & V & X & L & C & D & M\\
1 & 5 & 10 & 50 & 100 & 500 & 1.000\\
\end{tabular}
\end{center}

Πρόσθεταν από αριστερά προς τα δεξιά (II = 2, XVI = 16) και αφαιρούσαν όταν ένα μικρό σύμβολο προηγείται: IV = 4 ($5-1$), IX = 9, XL = 40, CM = 900. Τα σύμβολα που χρησιμοποιούμε σήμερα (0 ως 9) προέρχονται από τους Ινδούς, έφτασαν στην Ευρώπη μέσω των Αράβων και γι' αυτό λέγονται \textbf{ινδοαραβικά}. Εμφανίζονται πρώτη φορά στο έργο του Αλ-Χουαρίζμι (9ος αιώνας μ.Χ.).

Δεν είναι ανάγκη να έχουμε δέκα ψηφία. Ο υπολογιστής χρησιμοποιεί μόνο το 0 και το 1, στο \textbf{δυαδικό σύστημα}, όπου αντί για μονάδες, δεκάδες, εκατοντάδες έχουμε μονάδες, δυάδες, τετράδες, οκτάδες, δεκαεξάδες \dots

\begin{center}
\small\setlength{\tabcolsep}{4pt}
\begin{tabular}{l|cccccccccccc}
Δεκαδικό & 0 & 1 & 2 & 3 & 4 & 5 & 6 & 7 & 8 & 9 & 10 & 11\\
Δυαδικό & 0 & 1 & 10 & 11 & 100 & 101 & 110 & 111 & 1000 & 1001 & 1010 & 1011\\
\end{tabular}
\end{center}

Οι Πυθαγόρειοι πίστευαν ότι «τα πάντα είναι αριθμός». Έβαζαν τους αριθμούς σε σχήματα: οι \textbf{τρίγωνοι} αριθμοί $1,3,6,10,\dots$ φτιάχνουν τρίγωνα από τελείες και οι \textbf{τετράγωνοι} $1,4,9,16,\dots$ τετράγωνα. Οι τετράγωνοι είναι οι δεύτερες δυνάμεις.
\tmima{Λυμένα παραδείγματα}

\paradeigma{Οι κωδικοί της Λίνας}
\tixero Ο κωδικός έχει 4 ψηφία και κάθε ψηφίο έχει 10 επιλογές (0 ως 9).\par
\tithelo Πόσοι είναι οι δυνατοί κωδικοί, και πώς γράφεται σύντομα ο υπολογισμός;\par
\possindeo Για κάθε ψηφίο έχω 10 επιλογές, και οι επιλογές συνδυάζονται, άρα πολλαπλασιάζω: $10\cdot 10\cdot 10\cdot 10$. Είναι τέσσερα δεκάρια, δηλαδή $10^4$. Υπολογίζω βήμα βήμα: $10\cdot 10=100$, $100\cdot 10=1.000$, $1.000\cdot 10=10.000$.\par
\apantisi $10^4=10.000$ κωδικοί.\par
\elegxos Οι κωδικοί είναι οι αριθμοί από το 0000 ως το 9999. Από την παράγραφο 1.1: $9.999-0+1=10.000$. Βγήκε το ίδιο.

\paradeigma{Η παράσταση με πολλές πράξεις}
\tixero Η παράσταση $(2\cdot 5)^4+4\cdot(3+2)^2$.\par
\tithelo Την τιμή της.\par
\possindeo Πρώτα οι παρενθέσεις: $2\cdot 5=10$ και $3+2=5$. Έχω $10^4+4\cdot 5^2$. Μετά οι δυνάμεις: $10^4=10.000$ και $5^2=25$. Μετά ο πολλαπλασιασμός: $4\cdot 25=100$. Τέλος η πρόσθεση: $10.000+100=10.100$.\par
\apantisi $10.100$.\par
\elegxos Το $10^4$ έχει 4 μηδενικά, 10.000. Το άθροισμα με το 100 δίνει 10.100, όχι κάτι πιο μικρό από το 10.000. Σωστά.

\paradeigma{Ο αριθμός του Ηλία}
\tixero Ο Ηλίας γράφει ένα αριθμό σε ανάπτυγμα: $5\cdot 10^4+3\cdot 10^2+8$.\par
\tithelo Ποιος αριθμός είναι;\par
\possindeo $5\cdot 10^4=5\cdot 10.000=50.000$. $3\cdot 10^2=300$. Άρα $50.000+300+8=50.308$.\par
\apantisi $50.308$.\par
\elegxos Λείπουν οι χιλιάδες και οι δεκάδες, που σημαίνουν 0 στη θέση τους: 50.308. Τα μηδενικά έχουν μπει σωστά.

\tmima{Ασκήσεις}

\zoni{Το κατάλαβα}{εφαρμογή του κανόνα}

\ex Γράψε ως δύναμη: α) $5\cdot 5\cdot 5$ \quad β) $2\cdot 2\cdot 2\cdot 2\cdot 2$ \quad γ) $10\cdot 10\cdot 10\cdot 10$ \quad δ) $7\cdot 7$ \quad ε) $9$.

\ex Γράψε ως γινόμενο και υπολόγισε: α) $2^4$ \quad β) $3^3$ \quad γ) $5^2$ \quad δ) $10^5$ \quad ε) $1^{100}$ \quad στ) $6^1$.

\ex Ποια είναι η βάση και ποιος ο εκθέτης στα $4^5$, $12^2$, $1^{9}$ και $10^8$;

\ex Υπολόγισε το τετράγωνο και τον κύβο των αριθμών 2, 3, 4, 5 και 10. Φτιάξε πίνακα.

\ex Υπολόγισε την τιμή: α) $12+3\cdot 4$ \quad β) $(12+3)\cdot 4$ \quad γ) $24:4:2$ \quad δ) $100-5\cdot 4$ \quad ε) $2+3^2$ \quad στ) $(2+3)^2$.

\ex Βρες την τιμή: α) $(2\cdot 5)^4+4\cdot(3+2)^2$ \quad β) $(2+3)^3-8\cdot 3^2$ \quad γ) $3\cdot(4+2)^2-10\cdot 3$.

\ex Γράψε το ανάπτυγμα με δυνάμεις του 10: α) 7.604 \quad β) 53.280 \quad γ) 400.065.

\ex Γράψε με ψηφία τον αριθμό: α) $3\cdot 10^3+2\cdot 10^2+5$ \quad β) $9\cdot 10^5+4\cdot 10$ \quad γ) $10^6+10^2$.

\ex Σωστό ή λάθος; α) $2^3=6$ \quad β) $10^4=10.000$ \quad γ) $5^2=2\cdot 5$ \quad δ) $1^{57}=1$ \quad ε) $3+2\cdot 5=25$ \quad στ) $(1+1)^2=1^2+1^2$.

\zoni{Το χρησιμοποιώ}{προβλήματα από τη ζωή}

\ex \textbf{Το λουκέτο.} Ένα λουκέτο έχει τροχούς με ψηφία από το 0 ως το 9. Πόσους διαφορετικούς κωδικούς έχει αν έχει 3 τροχούς; Αν έχει 4; Αν έχει 5; Γράψε κάθε απάντηση και ως δύναμη του 10.

\ex \textbf{Οι διακόπτες.} Στην αποθήκη του θείου Λευτέρη ο πίνακας έχει διακόπτες που ο καθένας έχει δύο θέσεις: ανοιχτός ή κλειστός. Πόσους διαφορετικούς συνδυασμούς θέσεων έχουν 3 διακόπτες; Πόσους 5; Πόσους 10; Γράψε κάθε απάντηση και ως δύναμη του 2.

\ex \textbf{Ο κύβος του θείου.} Ο θείος Λευτέρης έχει δοχείο σε σχήμα κύβου με ακμή 4 δεκατόμετρα. Πόσα κυβικά δεκατόμετρα χωράει; (Θυμήσου ότι ο όγκος κύβου είναι $\alpha^3$.)

\ex \textbf{Ο κήπος.} Ο τετράγωνος κήπος της γιαγιάς Ρήνης έχει πλευρά 12 μ. Βρες το εμβαδό του. Αν διπλασιαστεί η πλευρά, ο κήπος γίνεται διπλάσιος σε εμβαδό; Υπολόγισε.

\ex \textbf{Τα σκακιστικά τετράγωνα.} Η σκακιέρα έχει 8 γραμμές και 8 στήλες. Πόσα τετράγωνα έχει; Γράψε τον υπολογισμό με δύναμη.

\ex \textbf{Ο λογαριασμός του Άρη.} Ο Άρης αγοράζει 3 μπάλες των 8 € και ένα φανελάκι των 12 €. Γράφει $3\cdot 8+12$ και βρίσκει 36. Η Μυρτώ γράφει $3\cdot(8+12)$. Ποια παράσταση ταιριάζει στο πρόβλημα και πόσο πληρώνει τελικά ο Άρης;

\ex \textbf{Τρεις απαντήσεις.} Τρεις φίλοι υπολογίζουν την παράσταση $4\cdot(7+7\cdot 9)+20$. Ο Άρης βρίσκει 335, η Ζωή 300 και ο Ηλίας 524. Ποιος έχει δίκιο; Βρες πού έκανε λάθος καθένας από τους άλλους δύο.

\zoni{Το σκέφτομαι}{μια ερώτηση που πάει πιο μακριά}

\ex Πόσο κάνει το $2^3\cdot 2^2$; Γράψε το ως γινόμενο παραγόντων (όλα τα δυάρια) και μέτρα. Ανακαλύπτεις έναν γρήγορο κανόνα για τον πολλαπλασιασμό δυνάμεων με την ίδια βάση;

\ex Συγκρίνε τα $2^5$ και $5^2$, και τα $2^4$ και $4^2$. Μετά βρες έναν αριθμό $\alpha$ διαφορετικό από το 2 για τον οποίο το $\alpha^2$ είναι ίσο με το $2\alpha$.

\ex Τοποθέτησε παρενθέσεις στην παράσταση $2+3\cdot 4+5$ με όσους τρόπους μπορείς και βρες όλες τις διαφορετικές τιμές που μπορεί να πάρει.

\ex Βάλε τα σύμβολα $+,-,\cdot$ (ή παρενθέσεις) ανάμεσα στα ψηφία $1\ \ 2\ \ 3\ \ 4$ ώστε το αποτέλεσμα να είναι 10. Βρες τουλάχιστον δύο τρόπους.

\ex Ο Ηλίας λέει: «Το $10^0$ πρέπει να είναι 0, γιατί δεν πολλαπλασιάζεις το 10 καμία φορά». Η Ζωή λέει: «Μάλλον είναι 1». Δες το μοτίβο $10^3=1.000$, $10^2=100$, $10^1=10$ και σκέψου τι πρέπει να συνεχίσει. Ποιος νομίζεις ότι έχει δίκιο; (Θα το μάθεις στο σχολείο σε λίγους μήνες.)

\zoni{Για το σπίτι}{εξάσκηση και ιστορία των αριθμών}

\ex Βρες τα τετράγωνα των 10, 20, 30, 40, 50, 60, 70, 80, 90 και τους κύβους των 10, 20, 30, 40, 50. Τι παρατηρείς στα μηδενικά;

\ex Υπολόγισε τις δυνάμεις $2^1, 2^2, 2^3, \dots, 2^{10}$ με τη σειρά, κάθε μία από την προηγούμενη.

\ex Γράψε πιο σύντομα: α) $\alpha+\alpha+\alpha$ \quad β) $\alpha\cdot\alpha\cdot\alpha$ \quad γ) $x+x+x+x$ \quad δ) $x\cdot x\cdot x\cdot x$ \quad ε) $2\cdot 2\cdot\alpha\cdot\alpha\cdot\alpha$.

\ex Βρες τις τιμές των παραστάσεων $(6+5)^2$ και $6^2+5^2$, και μετά των $(3+6)^2$ και $3^2+6^2$. Τι παρατηρείς;

\ex \textbf{Ρωμαϊκοί αριθμοί.} α) Γράψε με ρωμαϊκά τους 4, 9, 14, 40, 49, 90 και 1.994. β) Διάβασε τους XXXVII, CDXLIV και MMXXVI.

\ex \textbf{Αλφαβητικοί αριθμοί.} α) Γράψε με το αλφαβητικό σύστημα των Ιώνων τους 12, 25 και 831. β) Διάβασε τους ρκγ, σμ και τξε.

\ex \textbf{Δυαδικοί αριθμοί.} α) Γράψε στο δυαδικό σύστημα τους 11, 12 και 16. β) Διάβασε τους δυαδικούς 111, 1101 και 10101.

\ex \textbf{Οι αριθμοί των Πυθαγορείων.} Συνέχισε την ακολουθία των τριγώνων αριθμών $1,3,6,10$ για τρεις ακόμα όρους. Πόσο είναι ο δέκατος τρίγωνος αριθμός; Πόσο είναι ο δωδέκατος τετράγωνος;

\ex \textbf{Μαγικό τετράγωνο.} Βάλε τους αριθμούς από το 1 ως το 9, από μία φορά τον καθένα, ώστε κάθε γραμμή, κάθε στήλη και κάθε διαγώνιος να έχει άθροισμα 15. Έχουν μπει ήδη μερικοί:
\begin{center}
\renewcommand{\arraystretch}{1.6}\setlength{\tabcolsep}{12pt}
\begin{tabular}{|c|c|c|}\hline 2 & & 6\\\hline & 5 & \\\hline 4 & & 8\\\hline\end{tabular}
\end{center}
